\documentclass[12pt,a4paper]{article}
\usepackage[UTF8]{ctex}
\usepackage{geometry}
\usepackage{listings}
\usepackage{xcolor}
\usepackage{hyperref}
\usepackage{amsmath}
\usepackage{algorithm}
\usepackage{algpseudocode}

\geometry{margin=2.5cm}

\lstset{
    language=Python,
    basicstyle=\ttfamily\small,
    keywordstyle=\color{blue}\bfseries,
    commentstyle=\color{green!60!black},
    stringstyle=\color{red},
    numbers=left,
    numberstyle=\tiny\color{gray},
    stepnumber=1,
    numbersep=5pt,
    frame=single,
    breaklines=true,
    showstringspaces=false,
    tabsize=4
}

\title{Halo 系统代码文档：Schedulers 模块}
\author{代码分析文档}
\date{\today}

\begin{document}

\maketitle

\tableofcontents
\newpage

\section{概述}

\texttt{schedulers} 模块实现了多种调度策略，用于将工作流图中的 \texttt{Operator} 分配到多个 GPU 设备上执行。调度器的核心任务是：在满足依赖关系的前提下，生成每个设备上的执行序列（workflow），以优化整体执行时间和资源利用率。

\section{调度器接口}

所有调度器都遵循相同的接口：

\begin{lstlisting}
def schedule_xxx(
    device_cnt: int,           # 设备数量
    start_ops: List[Operator],  # 起始节点列表
    end_ops: List[Operator],   # 终止节点列表（部分调度器不使用）
    all_ops: List[Operator],   # 所有 OP 列表
    queries: List[Query],      # 查询列表
    **kwargs                   # 其他参数
) -> List[List[Dict]]:
    """
    返回 workflows，格式为：
    List[ per-device List[ {"command": "execute", "params": (op, query_ids)} ] ]
    """
\end{lstlisting}

\section{通用工具函数}

\subsection{\_topo\_order}

多个调度器共享的拓扑排序函数：

\begin{lstlisting}
def _topo_order(all_ops: List[Operator]) -> List[Operator]:
    """Kahn's algorithm for topological sort."""
    indeg = {op: len(op.input_ops) for op in all_ops}
    q = deque([op for op, d in indeg.items() if d == 0])
    topo = []
    while q:
        u = q.popleft()
        topo.append(u)
        for v in u.output_ops:
            indeg[v] -= 1
            if indeg[v] == 0:
                q.append(v)
    if len(topo) != len(all_ops):
        raise ValueError("Graph is not a DAG.")
    return topo
\end{lstlisting}

\begin{description}
    \item[算法] Kahn 算法（基于入度的拓扑排序）
    \item[时间复杂度] $O(V + E)$
    \item[功能] 返回一个满足依赖关系的 OP 执行顺序
    \item[用途] 确保在调度时，父节点总是先于子节点被分配
\end{description}

\section{调度策略详解}

\subsection{schedule\_rr：轮询调度}

\subsubsection{功能概述}

最简单的调度策略：按照拓扑顺序，将每个 OP 轮询分配到设备上，每个 OP 处理所有查询。

\subsubsection{实现代码}

\begin{lstlisting}
def schedule_rr(
    device_cnt: int,
    all_ops: List[Operator],
    queries: List[Query],
) -> List[List[Dict]]:
    """Round-Robin per op; each op handles ALL queries."""
    if device_cnt <= 0:
        raise RuntimeError("No devices available.")
    
    workflows: List[List[Dict]] = [[] for _ in range(device_cnt)]
    all_ids = [q.id for q in queries]
    topo = _topo_order(all_ops)
    
    d = 0
    for op in topo:
        workflows[d].append({"command": "execute", "params": (op, list(all_ids))})
        d = (d + 1) % device_cnt
    return workflows
\end{lstlisting}

\subsubsection{执行流程}

\begin{enumerate}
    \item 对所有 OP 进行拓扑排序
    \item 按顺序遍历每个 OP
    \item 将 OP 分配到设备 \texttt{d}，并将所有查询 ID 传递给该 OP
    \item 设备索引循环递增：\texttt{d = (d + 1) \% device\_cnt}
\end{enumerate}

\subsubsection{特点}

\begin{itemize}
    \item \textbf{优点}：实现简单，负载均衡（每个设备分配大致相同数量的 OP）
    \item \textbf{缺点}：
    \begin{itemize}
        \item 不考虑依赖关系的位置（父节点和子节点可能在不同设备）
        \item 不考虑模型切换成本（同一设备可能频繁切换模型）
        \item 不支持并行执行（同一时刻只有一个 OP 在执行）
    \end{itemize}
\end{itemize}

\subsubsection{适用场景}

\begin{itemize}
    \item 基准测试和对比
    \item 简单的工作流（OP 数量少，依赖关系简单）
    \item 所有 OP 使用相同模型的情况
\end{itemize}

\subsection{schedule\_by\_levels：层级调度}

\subsubsection{功能概述}

将 OP 按依赖层级分组，同一层级的 OP 可以并行执行，层内使用轮询分配。

\subsubsection{实现代码}

\begin{lstlisting}
def schedule_by_levels(
    device_cnt: int,
    all_ops: List[Operator],
    queries: List[Query],
) -> List[List[Dict]]:
    """Level-based RR: assign ops within the same topo level round-robin."""
    if device_cnt <= 0:
        raise RuntimeError("No devices available.")
    
    workflows: List[List[Dict]] = [[] for _ in range(device_cnt)]
    all_ids = [q.id for q in queries]
    topo = _topo_order(all_ops)
    
    # 计算每个 OP 的层级
    level = {}
    for op in topo:
        level[op] = 0 if not op.input_ops else max(level[p] for p in op.input_ops) + 1
    
    # 按层级分配
    d = 0
    for lv in sorted(set(level.values())):
        layer_ops = [op for op in topo if level[op] == lv]
        for op in layer_ops:
            workflows[d % device_cnt].append({
                "command": "execute",
                "params": (op, list(all_ids))
            })
            d += 1
    return workflows
\end{lstlisting}

\subsubsection{执行流程}

\begin{enumerate}
    \item 拓扑排序所有 OP
    \item \textbf{计算层级}：
    \begin{itemize}
        \item 起始节点（无输入）层级为 0
        \item 其他节点的层级 = 其所有父节点的最大层级 + 1
    \end{itemize}
    \item \textbf{按层级分组}：将同一层级的 OP 收集到一起
    \item \textbf{层级内轮询分配}：对每个层级内的 OP，按轮询方式分配到设备
\end{enumerate}

\subsubsection{层级计算示例}

假设工作流：\texttt{op0} $\rightarrow$ \texttt{op1} $\rightarrow$ \texttt{op2}

\begin{itemize}
    \item \texttt{op0}：层级 0（无输入）
    \item \texttt{op1}：层级 1（依赖 op0）
    \item \texttt{op2}：层级 2（依赖 op1）
\end{itemize}

\subsubsection{特点}

\begin{itemize}
    \item \textbf{优点}：
    \begin{itemize}
        \item 同一层级的 OP 可以并行执行（如果设备数足够）
        \item 保证依赖关系：低层级总是先于高层级执行
    \end{itemize}
    \item \textbf{缺点}：
    \begin{itemize}
        \item 仍然不考虑模型切换成本
        \item 不考虑父节点和子节点的设备位置（可能跨设备传输数据）
    \end{itemize}
\end{itemize}

\subsection{schedule\_colocate\_parents：父节点共置调度}

\subsubsection{功能概述}

优先将 OP 分配到其父节点所在的设备，以减少跨设备数据传输，提高缓存局部性。

\subsubsection{实现代码}

\begin{lstlisting}
def schedule_colocate_parents(
    device_cnt: int,
    all_ops: List[Operator],
    queries: List[Query],
) -> List[List[Dict]]:
    """
    Prefer assigning an op to any parent's device (locality).
    Each op uses ALL queries.
    """
    if device_cnt <= 0:
        raise RuntimeError("No devices available.")
    
    workflows: List[List[Dict]] = [[] for _ in range(device_cnt)]
    all_ids = [q.id for q in queries]
    topo = _topo_order(all_ops)
    
    op2dev: Dict[Operator, int] = {}  # OP -> 设备映射
    rr = 0  # 轮询计数器（用于无父节点的情况）
    
    for op in topo:
        chosen = None
        # 尝试找到父节点所在的设备
        for p in op.input_ops:
            if p in op2dev:
                chosen = op2dev[p]
                break
        # 如果没有父节点或父节点未分配，使用轮询
        if chosen is None:
            chosen = rr
            rr = (rr + 1) % device_cnt
        
        op2dev[op] = chosen
        workflows[chosen].append({
            "command": "execute",
            "params": (op, list(all_ids))
        })
    return workflows
\end{lstlisting}

\subsubsection{执行流程}

\begin{enumerate}
    \item 拓扑排序所有 OP
    \item 按顺序遍历每个 OP：
    \begin{enumerate}
        \item 检查该 OP 的所有父节点
        \item 如果某个父节点已经分配到设备，则将该 OP 分配到同一设备
        \item 如果没有父节点或父节点都未分配，使用轮询选择设备
    \end{enumerate}
    \item 记录 OP 到设备的映射
\end{enumerate}

\subsubsection{特点}

\begin{itemize}
    \item \textbf{优点}：
    \begin{itemize}
        \item 提高缓存局部性：父节点和子节点在同一设备，可以减少数据传输
        \item 如果使用 KV cache，可以更好地利用缓存
    \end{itemize}
    \item \textbf{缺点}：
    \begin{itemize}
        \item 可能导致负载不均衡（某些设备承担更多工作）
        \item 不考虑模型切换成本
    \end{itemize}
\end{itemize}

\subsection{schedule\_dp：数据并行调度}

\subsubsection{功能概述}

将查询 ID 固定分片到各设备，每个 OP 在所有设备上并行执行，但每个设备只处理分配给它的查询子集。

\subsubsection{实现代码}

\begin{lstlisting}
def schedule_dp(
    device_cnt: int,
    all_ops: List[Operator],
    queries: List[Query],
) -> List[List[Dict]]:
    """
    Sticky data-parallel with fixed shards:
      - Split query ids into D shards.
      - For each op in topo order, submit (op, shard_d) to device d.
    """
    if device_cnt <= 0:
        raise RuntimeError("No devices available.")
    
    workflows: List[List[Dict]] = [[] for _ in range(device_cnt)]
    topo = _topo_order(all_ops)
    D = max(1, device_cnt)
    all_ids = [q.id for q in queries]
    
    # 将查询 ID 分片
    shards = []
    N = len(all_ids)
    for d in range(device_cnt):
        start = (N * d) // D
        end = (N * (d + 1)) // D
        shard = all_ids[start:end]
        shards.append(shard)
    
    # 每个 OP 在所有设备上执行，但只处理对应的查询分片
    for op in topo:
        for d, shard in enumerate(shards):
            if not shard:
                continue
            workflows[d].append({
                "command": "execute",
                "params": (op, list(shard))
            })
    return workflows
\end{lstlisting}

\subsubsection{执行流程}

\begin{enumerate}
    \item 拓扑排序所有 OP
    \item \textbf{查询分片}：
    \begin{itemize}
        \item 将查询 ID 列表均匀分成 \texttt{device\_cnt} 个分片
        \item 每个设备对应一个分片
    \end{itemize}
    \item \textbf{为每个 OP 创建任务}：
    \begin{itemize}
        \item 对每个 OP，在所有设备上都创建一个执行任务
        \item 每个设备只处理分配给它的查询分片
    \end{itemize}
\end{enumerate}

\subsubsection{查询分片示例}

假设有 10 个查询（ID: 0-9），3 个设备：

\begin{itemize}
    \item 设备 0：查询 [0, 1, 2, 3]
    \item 设备 1：查询 [4, 5, 6]
    \item 设备 2：查询 [7, 8, 9]
\end{itemize}

每个 OP 都会在这 3 个设备上并行执行，但每个设备只处理自己的查询分片。

\subsubsection{特点}

\begin{itemize}
    \item \textbf{优点}：
    \begin{itemize}
        \item 真正的并行执行：同一 OP 在多个设备上同时运行
        \item 负载均衡：每个设备处理相同数量的查询
        \item 适合大批量查询的场景
    \end{itemize}
    \item \textbf{缺点}：
    \begin{itemize}
        \item 每个 OP 都需要在所有设备上加载模型（如果模型不同）
        \item 可能导致显存占用较高
    \end{itemize}
\end{itemize}

\subsection{schedule\_heuristic：启发式调度}

\subsubsection{功能概述}

最复杂的调度策略之一，结合了多种优化技术：
\begin{itemize}
    \item 依赖感知的前沿扩展
    \item 基于 \texttt{max\_distance} 的优先级排序
    \item 数据并行复制（当设备数多于就绪 OP 数时）
    \item 父节点设备优先分配（缓存局部性）
\end{itemize}

\subsubsection{实现代码（关键部分）}

\begin{lstlisting}
def schedule_heuristic(
    device_cnt: int,
    start_ops: List[Operator],
    end_ops: List[Operator],
    all_ops: List[Operator],
    queries: List[Query],
    dp_threshold: int = 2,
) -> List[List[Dict]]:
    """
    Heuristic scheduler:
      - Grow frontier respecting dependencies.
      - If ops > devices: pick by longest distance to end (desc).
      - If ops < devices: try to resume paused; else DP duplicate selected ops.
      - Prefer assigning to parent's device (cache locality).
    """
    if device_cnt <= 0:
        raise RuntimeError("No devices available.")
    
    _compute_longest_distances(all_ops, end_ops)
    
    workflows: List[List[Dict]] = [[] for _ in range(device_cnt)]
    requests_cnt = len(queries)
    last_ops: List[Operator] = []  # 上一轮执行的 OP
    paused_ops: List[Operator] = []  # 暂停的 OP（依赖未满足）
    device_history: Dict[Operator, int] = {}  # OP -> 设备映射
    
    while True:
        # 1) 计算前沿（Frontier）
        if not last_ops:
            new_ops = list(start_ops)
        else:
            new_ops = []
            for op in last_ops:
                new_ops.extend(op.output_ops)
            new_ops = list(set(new_ops))
        
        # 2) 检查依赖是否满足
        valid = []
        for op in new_ops:
            if all(inp in device_history for inp in op.input_ops):
                valid.append(op)
                if op in paused_ops:
                    paused_ops.remove(op)
            else:
                if op not in paused_ops:
                    paused_ops.append(op)
        new_ops = valid
        
        # 3) 尝试恢复暂停的 OP
        if not new_ops:
            resumed = []
            for op in paused_ops[:]:
                if all(inp in device_history for inp in op.input_ops):
                    resumed.append(op)
                    paused_ops.remove(op)
            new_ops.extend(resumed)
        
        if not new_ops:
            break
        
        # 4) 适配设备数量
        if len(new_ops) > device_cnt:
            # 按 max_distance 降序排序，选择前 device_cnt 个
            new_ops.sort(key=lambda x: getattr(x, "max_distance", 0), reverse=True)
            extra = new_ops[device_cnt:]
            new_ops = new_ops[:device_cnt]
            for op in extra:
                if op not in paused_ops:
                    paused_ops.append(op)
        
        elif len(new_ops) < device_cnt:
            # 尝试恢复更多暂停的 OP
            for op in paused_ops[:]:
                if all(inp in device_history for inp in op.input_ops):
                    new_ops.append(op)
                    paused_ops.remove(op)
                    if len(new_ops) == device_cnt:
                        break
            
            # 如果仍然不足，且查询数足够，进行数据并行复制
            if len(new_ops) < device_cnt and requests_cnt > dp_threshold and len(new_ops) > 0:
                required = device_cnt - len(new_ops)
                duplicates: List[Operator] = []
                for i in range(required):
                    original = new_ops[i % len(new_ops)]
                    # 创建复制节点...
                    # （详细代码见源码）
                    duplicates.append(dup)
                new_ops.extend(duplicates)
        
        # 5) 设备分配（优先父节点设备）
        available = set(range(device_cnt))
        for op in new_ops:
            assigned = False
            for inp in op.input_ops:
                if inp in last_ops:
                    d = device_history[inp]
                    if d in available:
                        device_history[op] = d
                        available.remove(d)
                        assigned = True
                        break
            if not assigned:
                device_history[op] = available.pop()
        
        # 6) 生成执行任务
        for op in new_ops:
            req_ids = _partition_query_ids(op, queries)
            workflows[device_history[op]].append({
                "command": "execute",
                "params": (op, req_ids)
            })
        
        last_ops = new_ops
    
    return workflows
\end{lstlisting}

\subsubsection{关键步骤详解}

\begin{enumerate}
    \item \textbf{前沿计算}：
    \begin{itemize}
        \item 第一轮：从 \texttt{start_ops} 开始
        \item 后续轮：从前一轮执行的 OP 的后继节点中收集
    \end{itemize}
    
    \item \textbf{依赖检查}：
    \begin{itemize}
        \item 只有所有父节点都已分配设备的 OP 才能执行
        \item 不满足的 OP 加入 \texttt{paused_ops}
    \end{itemize}
    
    \item \textbf{优先级排序}：
    \begin{itemize}
        \item 当就绪 OP 数多于设备数时，按 \texttt{max_distance} 降序排序
        \item 优先执行距离终点更远的 OP（尽早开始长链）
    \end{itemize}
    
    \item \textbf{数据并行复制}：
    \begin{itemize}
        \item 当就绪 OP 数少于设备数，且查询数足够时
        \item 创建 OP 的副本，分配到不同设备
        \item 查询 ID 被分片到各副本
    \end{itemize}
    
    \item \textbf{设备分配策略}：
    \begin{itemize}
        \item 优先分配到父节点所在的设备（缓存局部性）
        \item 如果父节点不在上一轮，则使用可用设备
    \end{itemize}
\end{enumerate}

\subsubsection{\_partition\_query\_ids 函数}

\begin{lstlisting}
def _partition_query_ids(op: Operator, queries: List[Query], return_full: bool = False):
    """Partition queries among duplicates; otherwise return all query ids."""
    if getattr(op, "data_parallel", False):
        dup_index, total_dup = op.duplicate_info[0], op.duplicate_info[1] + 1
        if total_dup > 0:
            total = len(queries)
            per = total // total_dup
            start = per * dup_index
            end = per * (dup_index + 1) if dup_index < total_dup else total
            queries = queries[start:end]
    if return_full:
        return queries
    return [x.id for x in queries]
\end{lstlisting}

\begin{description}
    \item[功能] 根据 OP 是否数据并行，返回对应的查询 ID 列表
    \item[逻辑]
    \begin{itemize}
        \item 如果 OP 是数据并行的副本，则只返回分配给该副本的查询 ID 分片
        \item 否则返回所有查询 ID
    \end{itemize}
\end{description}

\subsection{schedule\_search：Beam-Search 调度}

\subsubsection{功能概述}

最复杂的调度策略，使用 Beam-Search 算法搜索近似最优的设备分配方案，考虑模型切换成本、负载均衡等因素。

\subsubsection{核心思想}

\begin{itemize}
    \item 维护一个 beam（候选方案集合）
    \item 每轮扩展：为当前就绪的 OP 生成所有可能的设备分配方案
    \item 评估每个方案的代价（makespan + 模型切换成本）
    \item 保留最优的 K 个方案（beam width）
    \item 重复直到所有 OP 都被分配
\end{itemize}

\subsubsection{关键数据结构}

\begin{lstlisting}
@dataclass(frozen=True)
class Placement:
    """单个阶段的设备分配：哪个 OP 在哪个设备上运行"""
    op: Operator
    rep_idx: int = 0      # 副本索引（如果复制）
    rep_total: int = 1    # 总副本数

@dataclass
class BeamState:
    """部分分配状态"""
    seqs: List[List[Placement]]  # 每个设备的执行序列
    assigned: Set[Operator]      # 已分配的 OP 集合
    cost_accum: List[float]      # 每个设备的累计代价
    score: float = 0.0            # 全局评分
\end{lstlisting}

\subsubsection{代价模型}

\begin{lstlisting}
PER_TOKEN_COST = 1.0          # 每个生成 token 的基础代价
MODEL_SWITCH_PENALTY = 300.0  # 模型切换惩罚
SAME_MODEL_BONUS = -80.0      # 相同模型连续执行的奖励

def _estimate_device_inc_cost(
    last_op: Optional[Operator],
    new_op: Operator,
    num_queries: int,
) -> float:
    """计算在设备上追加新 OP 的增量代价"""
    cost = _op_demand(new_op, num_queries)  # 基础计算需求
    
    if last_op is None:
        return cost  # 第一个任务
    
    last_model = last_op.model_config.model_name
    new_model = new_op.model_config.model_name
    if last_model != new_model:
        cost += MODEL_SWITCH_PENALTY  # 模型切换惩罚
    else:
        cost += SAME_MODEL_BONUS      # 相同模型奖励
    
    return cost
\end{lstlisting}

\subsubsection{执行流程（简化）}

\begin{enumerate}
    \item 初始化 beam，包含一个空状态
    \item 循环直到所有 OP 被分配：
    \begin{enumerate}
        \item 对 beam 中的每个状态：
        \begin{enumerate}
            \item 找出所有就绪的 OP（依赖已满足）
            \item 选择最多 \texttt{device\_cnt} 个就绪 OP
            \item 如果不足，考虑复制最重的 OP
            \item 生成设备分配方案（考虑父节点设备优先）
            \item 计算增量代价和评分
            \item 生成新的候选状态
        \end{enumerate}
        \item 保留评分最优的 K 个状态（beam width）
    \end{enumerate}
    \item 返回最优状态的 workflows
\end{enumerate}

\subsubsection{评分函数}

\begin{lstlisting}
def _score_state(state: BeamState, all_ops: List[Operator], num_queries: int, device_cnt: int) -> float:
    """
    评分 = 当前最大 makespan + 剩余工作量的乐观下界
    """
    current = max(state.cost_accum) if state.cost_accum else 0.0
    
    remaining_ops = [op for op in all_ops if op not in state.assigned]
    remaining_demand = sum(_op_demand(op, num_queries) for op in remaining_ops)
    optimistic_lb = remaining_demand / max(1, device_cnt)
    
    return current + optimistic_lb
\end{lstlisting}

\subsubsection{特点}

\begin{itemize}
    \item \textbf{优点}：
    \begin{itemize}
        \item 考虑模型切换成本，减少不必要的模型重载
        \item 考虑负载均衡，尽量使各设备完成时间接近
        \item 支持 OP 复制和数据并行
    \end{itemize}
    \item \textbf{缺点}：
    \begin{itemize}
        \item 计算开销较大（需要维护 beam 和评估多个候选）
        \item 参数调优复杂（beam width、代价权重等）
    \end{itemize}
\end{itemize}

\section{调度器对比}

\begin{table}[h]
\centering
\begin{tabular}{|l|c|c|c|c|}
\hline
调度器 & 并行执行 & 模型切换优化 & 缓存局部性 & 复杂度 \\
\hline
RR & 否 & 否 & 否 & 低 \\
\hline
By Levels & 是（层级内） & 否 & 否 & 低 \\
\hline
Colocate Parents & 否 & 否 & 是 & 中 \\
\hline
Data Parallel & 是 & 否 & 否 & 中 \\
\hline
Heuristic & 是 & 部分 & 是 & 高 \\
\hline
Beam-Search & 是 & 是 & 是 & 很高 \\
\hline
\end{tabular}
\caption{调度器特性对比}
\end{table}

\section{总结}

\texttt{schedulers} 模块提供了从简单到复杂的多种调度策略：

\begin{itemize}
    \item \textbf{简单策略}（RR、By Levels）：适合快速原型和基准测试
    \item \textbf{中等策略}（Colocate Parents、Data Parallel）：针对特定优化目标
    \item \textbf{复杂策略}（Heuristic、Beam-Search）：综合考虑多种因素，追求整体最优
\end{itemize}

选择合适的调度器需要根据：
\begin{itemize}
    \item 工作流复杂度
    \item 设备数量和查询数量
    \item 模型切换成本的重要性
    \item 对执行时间的敏感度
\end{itemize}

\end{document}
